\chapter{Allergy Skin Testing}

\section*{Overview}
Allergy skin testing is an in-vivo diagnostic procedure that observes the skin's reaction to small amounts of suspected allergens. It helps identify triggers for allergic symptoms. The exact approach, setting, and preparation depend on the indication, anatomy, and the patient's health.

\section*{Why It Is Done}
It may evaluate allergic rhinitis, asthma, eczema, food allergy, medication allergy, or reactions to insect venom. The results are interpreted together with the person's history because a positive reaction does not always mean that an allergen causes symptoms.

\section*{How It Is Performed}
For a skin-prick test, drops of allergen extracts are placed on the skin, usually the forearm or back, and the surface is lightly scratched or pricked. For selected situations, a small amount may be injected into the skin for intradermal testing. The clinician measures raised, itchy wheals after a short observation period.

\section*{Preparation}
Some antihistamines and other medicines can interfere with results and may need to be stopped only as directed by the clinician. Tell the allergy specialist about medicines, asthma control, pregnancy, and previous severe reactions. Do not stop essential medicines without advice.

\section*{Risks and Results}
Itching and local swelling are common and usually temporary. Rarely, a more serious allergic reaction can occur, so testing should be performed where appropriate treatment is available. Results help guide avoidance, medication choices, or allergen immunotherapy; negative testing does not exclude every type of allergy.

\section*{Contraindications and Precautions}
There are usually no major absolute contraindications. Testing may be delayed over active skin disease or during acute illness. The clinician reviews asthma control, medicines that may affect treatment of a reaction, pregnancy, and the need for observation; do not stop medicines without instructions.

\section*{When to Seek Medical Care}
Follow the procedure team's specific instructions. Seek urgent medical assessment for severe or worsening pain, uncontrolled bleeding, fever or signs of infection, trouble breathing, chest pain, fainting, new weakness or confusion, inability to urinate, or any rapidly worsening symptom. The appropriate warning signs depend on the procedure and the patient's medical condition.

\section*{References}
\begin{enumerate}
  \item MedlinePlus. Allergy skin tests. U.S. National Library of Medicine; 2025.
  \item Current Medical Diagnosis and Treatment. McGraw-Hill; 2025.
\end{enumerate}
\clearpage
